\section{Discussion and attribution}
The modern experiments establish that detector operating point affects
downstream MOT, that AC--MOT v3 changes the declared high-resolution baseline on
two tracker hosts, and that operating-point choices shape measured compute.
They do not isolate a separate accuracy contribution from the exact
scene-conditioned switching schedule.

A matched-static control achieved comparable HOTA (44.530 versus 44.523 for
ByteTrack; 48.269 versus 48.095 for OATrack), and shuffled timing did not
underperform the causal schedule on development ($\Delta$HOTA = 2.617228 versus
2.446957). The present experiments therefore support operating-point
calibration and observation-set shaping, but do not isolate an additional
accuracy benefit from exact switching timing. Full attribution-control details
are provided in Online Resource~1.

The systems interpretation is still meaningful. A finite policy can allocate
detector compute online, react to changing scene statistics, and expose a
reproducible operating-point mixture without retraining detector weights. This
may be useful under an average detector-time budget even when a matched static
point has similar HOTA. It is a deployment interpretation, not evidence that
dynamic timing is causally superior.

Tracker dependence reinforces the need to keep hosts separate. ByteTrack starts
from a weaker modern baseline and shows the larger VisDrone change; OATrack
starts stronger and shows a smaller HOTA increment and essentially flat UAVDT
HOTA transfer. The detector-side policy changes the observations available to
each association algorithm, so a pooled ``AC--MOT gain'' would be misleading.
The negative searches and the documented Trial 7 selection deviation further
bound the conclusion: this paper reports the frozen record and does not rerun
the search.

The strongest supported conclusion is that AC--MOT is an auditable,
training-free mechanism for shaping the detector observation stream and
computational operating point of aerial MOT. It can improve a declared
baseline and transfer with tracker-dependent outcomes, but the current evidence
does not establish a separate, consistent timing-specific accuracy benefit.
